%%%%%%%%%%%%%%%%%%%%%%%%%%%%%%%%%%%%%%%%%%%%%%%%%%%%%%%%%%%%%%%%%%%%%%
% Lớp 11 | Dãy số. Cấp số cộng. Cấp số nhân | Cấp số cộng
% Dạng 1D2?2-4: Tìm số hạng cụ thể trong cấp số cộng
% Số câu: 92  (sinh tự động bởi tach_cau.py; từ nay thầy sửa trực tiếp file này)
%%%%%%%%%%%%%%%%%%%%%%%%%%%%%%%%%%%%%%%%%%%%%%%%%%%%%%%%%%%%%%%%%%%%%%

%%%%%%%%%%%%%%%% TRẮC NGHIỆM 4 PHƯƠNG ÁN %%%%%%%%%%%%%%%%

% nguồn: Ok_Toan11_TracNghiem_Lan2
\begin{ex}%[1D2N2-4]%[ma:90ede2b9]
	Tìm số hạng thứ $11$ của cấp số cộng có số hạng đầu bằng $3$ và công  sai $d=-2$.
	\choice 
		 {$-21$}
		 {$23$}
		 {$-19$}
		{\True $-17$} 
	\loigiai{ 
		Áp dụng công thức số hạng tổng quát của cấp số cộng ta có $u_{11} = u_1 +  10\cdot d = 3 + 10\cdot (-2) = -17$.
	} 
\end{ex}

% nguồn: Ok_Toan11_TracNghiem_Lan2
\begin{ex}%[1D2N2-4]%[ma:27d6c0a8]
	 Cho cấp số cộng $(u_n)$ có số hạng đầu $u_1=2$ và công sai $d=4$. Giá trị $u_{2019}$ bằng
	\choice 
		 {\True $8074$}
		 {$4074$}
		 {$8078$}
		{$4078$} 
	 \loigiai{
		Áp dụng công thức của số  hạng tổng quát $u_n=u_1+(n-1)d=2+2018\cdot4=8074$.
	} 
\end{ex}

% nguồn: Ok_Toan11_TracNghiem_Lan2
\begin{ex}%[1D2N2-4]%[ma:d11b71e0]
	Cho cấp số cộng $(u_n)$ có $u_1=11$ và  công sai $d=4$. Hãy tính $u_{99}$.
	 \choice
		 {$401$}
		{\True $403$} 
		{$402$} 
		 {$404$}
	 \loigiai{
		Ta có $u_{99}=u_1+98d=11+98\cdot4=403$. 
	} 
\end{ex}

% nguồn: Ok_Toan11_TracNghiem_Lan2
\begin{ex}%[1D2N2-4]%[ma:dac70bc5]
	Cho  cấp số cộng $(u_n)$ có số hạng đầu $u_1=2$ và công sai $d=4$. Giá trị $u_{2019}$ bằng
	\choice 
		 {$4\,074$}
		 {$8\,078$}
		 {$4\,078$}
		 {\True $8\,074$}
	\loigiai{ 
		Áp dụng công thức của số hạng tổng quát $u_n = u_1 + (n-1)\cdot d \Rightarrow u_{2019}=  2 + 2\,018\cdot 4 = 8\,074$.
	 }
\end{ex}

% nguồn: Ok_Toan11_TracNghiem_Lan2
\begin{ex}%[1D2N2-4]%[ma:fdcc7a2b]
	Cho  cấp số cộng $(u_n)$ có số hạng đầu $u_1=2$ và công sai $d=5$. Giá trị $u_4$ bằng
	 \choice
		{$22$} 
		 {$12$}
		{$250$} 
		 {\True $17$}
	 \loigiai{
		Ta có $u_4 = u_1 + 3\cdot d =  2 + 3\cdot 5 = 17$.
	 }
\end{ex}

% nguồn: Ok_Toan11_TracNghiem_Lan2
\begin{ex}%[1D2N2-4]%[ma:06baa60d]
	Một cấp số cộng $(u_n)$ có $u_{13}=8$ và $d=-3$.  Tìm số hạng thứ ba của cấp số cộng $(u_n)$.
	 \choice
		{$50$} 
		{$28$} 
		{\True $38$} 
		 {$44$}
	 \loigiai{
		Ta có $u_{13}=u_1+12d  \Rightarrow 8=u_1+12\cdot(-3) \Rightarrow u_1=44\Rightarrow u_3=u_1+2d=44-6=38$.
	 }
\end{ex}

% nguồn: Ok_Toan11_TracNghiem_Lan2
\begin{ex}%[1D2N2-4]%[ma:ee1a1b16]
	Cho cấp số cộng $(u_n)$ biết $u_2=-1$;  $u_4=7$. Tìm $u_3$.
	 \choice
		{$4$} 
		 {$10$}
		 {$8$}
		{\True  $3$}
	\loigiai{ 
		Áp dụng tính chất  của các số hạng trong dãy cấp số cộng, ta có $u_3=\dfrac{u_2+u_4}{2}=\dfrac{-1+7}{2}=3$.
	 }
\end{ex}

% nguồn: Ok_Toan11_TracNghiem_Lan2
\begin{ex}%[1D2N2-4]%[ma:052ae601]
	Viết ba số hạng xen giữa $2$ và $22$ để được một cấp số cộng có năm số  hạng. Ba số hạng đó lần lượt là
	\choice 
		{\True $7;  \ 12; \ 17$}
		{$6;  \ 10; \ 14$}
		{$8; \ 13;  \ 18$}
		{$6;  \ 12; \ 18$}
	\loigiai{ 
		Ta có $u_1=2$, $u_5=22$. Áp dụng công thức $d = \dfrac{u_n-u_1}{n-1}$ ta  được
		\[ 
		d=\dfrac{u_5-u_1}{4}=\dfrac{22-2}{4}=5. 
		\] 
		Suy ra 
		 \[
		\begin{aligned}
			 u_2&=u_1+d=7; \\
			 u_3&=u_1+2d=12; \\
			 u_4&=u_1+3d=17.
		\end{aligned}
		 \]
		Vậy ba số hạng cần  tìm là $7; \ 12; \ 17$.
	} 
\end{ex}

% nguồn: Ok_Toan11_TracNghiem_Lan2
\begin{ex}%[1D2N2-4]%[ma:db31d1d9]
	Cho cấp số cộng $(u_n)$ với $u_1=8$ và công sai $d=3$.  Giá trị của $u_2$ bằng
	\choice 
		 {$\dfrac{8}{3}$}
		{$24$} 
		 {$5$}
		{\True $11$} 
	\loigiai{Ta  có $u_2=u_1+d=8+3=11$.
	 }
\end{ex}

% nguồn: Ok_Toan11_TracNghiem_Lan2
\begin{ex}%[1D2N2-4]%[ma:a2ff8ee7]
	Cho cấp số cộng $(u_n)$ với $u_1=2022$ và công sai $d=7$. Giá trị  của $u_6$ bằng
	 \choice
		{$2043$} 
		 {$2064$}
		 {$2050$}
		{\True  $2057$}
	\loigiai{Ta có $u_6=u_1+5d=2022+5 \cdot 7 =  2057$.
	} 
\end{ex}

% nguồn: Ok_Toan11_TracNghiem_Lan2
\begin{ex}%[1D2N2-4]%[ma:00b83225]
	Cho cấp số cộng $(u_n)$ có số hạng đầu $u_1=-2$ và công sai $d=-7$. Giá trị  $u_6$ bằng
	\choice 
		{$37$} 
		{\True  $-37$}
		 {$-33$}
		{$33$} 
	\loigiai{ 
		Ta có $u_6=u_1+5d=-2-35=-37$. 
	} 
\end{ex}

% nguồn: Ok_Toan11_TracNghiem_Lan2
\begin{ex}%[1D2N2-4]%[ma:cc8c6dbc]
	Cho cấp số cộng $(u_n)$ với số hạng đầu tiên $u_1=2$ và công  sai $d=2$. Tìm $u_{2018}$?
	 \choice
		 {$u_{2018}=2^{2018}$}
		{$u_{2018}=2^{2017}$} 
		{\True $u_{2018}=4036$} 
		 {$u_{2018}=4038$}
	 \loigiai{
		 Ta có $u_n=u_1+(n-1)d \Rightarrow u_{2018}=2+(2018-1)\cdot2=4036$.
	} 
\end{ex}

% nguồn: Ok_Toan11_TracNghiem_Lan2
\begin{ex}%[1D2N2-4]%[ma:9859cbc6]
	Cho cấp số cộng  có $u_1=-3$, $d=4$. Chọn khẳng định đúng trong các khẳng định sau?
	 \choice
		 {$u_5=15$}
		{$u_4=8$} 
		{\True $u_3=5$} 
		 {$u_2=2$}
	 \loigiai{
		Ta có  $u_3=u_1+2d=-3+2\cdot4=5$.
	 }
\end{ex}

% nguồn: Ok_Toan11_TracNghiem_Lan2
\begin{ex}%[1D2N2-4]%[ma:979184ae]
	Cho cấp số cộng $1$, $4$, $7$, $\ldots$. Số  hạng thứ $100$ của cấp số cộng là
	\choice 
		 {$297$}
		{$301$} 
		{$295$} 
		 {\True $298$}
	\loigiai{ 
		Cấp số cộng $1$, $4$,  $7$, $\ldots$. có số hạng đầu $u_1=1$ và công sai $d=3$.\\
		Số hạng thứ $100$ của cấp số cộng là $u_{100}=u_1+99\cdot d=1+99\cdot3=298$. 
	 }
\end{ex}

% nguồn: Ok_Toan11_TracNghiem_Lan2
\begin{ex}%[1D2N2-4]%[ma:57037f61]
	Cho cấp số cộng $(u_n)$, biết $u_1=3$ và $u_2=-1$. Khi đó $u_3$ bằng 
	\choice 
		 {$u_3=4$}
		{$u_3=7$} 
		{$u_3=2$} 
		 {\True $u_3=-5$}
	 \loigiai{
		Ta có  $(u_n)$ là cấp số cộng nên $2u_2=u_1+u_3$ suy ra $u_3=2u_2-u_1=-5$.
	} 
\end{ex}

% nguồn: Ok_Toan11_TracNghiem_Lan2
\begin{ex}%[1D2N2-4]%[ma:590e8cfd]
	Cho cấp số cộng $(u_n)$ với $u_1=10$;  $u_2=13$. Giá trị của $u_4$ là
	 \choice
		{$u_4=20$} 
		{\True  $u_4=19$}
		{$u_4=16$} 
		{$u_4=18$} 
	 \loigiai{
		 $u_1=10; u_2=13 \Rightarrow d=3 \Rightarrow u_4=u_1+3\cdot d=10+3\cdot 3=19$.
	 }
\end{ex}

% nguồn: Ok_Toan11_TracNghiem_Lan2
\begin{ex}%[1D2N2-4]%[ma:138d8769]
	Cho cấp số cộng $(u_n)$ với $u_1=9$ và công sai $d=2$. Giá  trị của $u_2$ bằng
	\choice 
		{\True$11$} 
		{$\dfrac{9}{2}$} 
		 {$18$}
		 {$7$}
	 \loigiai{Ta có $u_2=u_1+d=9+2=11$.
	 }
\end{ex}

% nguồn: Ok_TracNghiemLop11_Lan1
\begin{ex}%[1D2N2-4]%[ma:cfadeaee]
	Cho cấp số cộng có số hạng đầu $u_1=-\dfrac{1}{2}$, công sai $d=\dfrac{1}{2}$. Năm số hạng liên tiếp đầu tiên  của cấp số này là
	 \choice
		 {$-\dfrac{1}{2};0;1;\dfrac{1}{2};1$}
		{$-\dfrac{1}{2};0;\dfrac{1}{2};0;\dfrac{1}{2}$} 
		{$\dfrac{1}{2};1;\dfrac{3}{2};2;\dfrac{5}{2}$} 
		{\True $-\dfrac{1}{2};0;\dfrac{1}{2};1;\dfrac{3}{2}$} 
	 \loigiai{
		Ta dùng công thức tổng quát $u_n=u_1+\left(n-1\right)d=-\dfrac{1}{2}+\left(n-1\right)\dfrac{1}{2}=-1+\dfrac{n}{2}$, hoặc $u_{n+1}=u_n+d=u_n+\dfrac{1}{2}$ để tính các số hạng của một cấp  số cộng.\\
		 Ta có $u_1=-\dfrac{1}{2}; d=\dfrac{1}{2} \Rightarrow \heva{&{u_1=-\dfrac{1}{2}} \\
		&{u_2=u_1+d=0} \\ 
		 &{u_3-u_2+d=\dfrac{1}{2}} \\
		&{u_4=u_3+d=1}  \\
		 &{u_5=u_4+d=\dfrac{3}{2}.}
		}$ 
		\\ \textbf{Nhận xét}:  Dùng chức năng "lặp" của MTCT để tính\\
		\begin{itemize}
			\item Nhập: $X=X+\dfrac{1}{2}$ (nhập $X=X+d$). 
			\item Bấm  CALC nhập $-\dfrac{1}{2}$ (nhập $u_1$).
			\item Để tính  $5$ số hạng đầu ta bấm dấu "=" liên tiếp để ra kết quả $4$ lần nữa!
		\end{itemize}
	} 
\end{ex}

% nguồn: Ok_TracNghiemLop11_Lan1
\begin{ex}%[1D2N2-4]%[ma:c4c3f670]
	Cho cấp số cộng $(u_n)$ có $u_1=-1$, công sai $d=3$. Số hạng thứ  $675$ của cấp số cộng đã cho là
	\choice 
		 {$2024$}
		 {$-2022$}
		 {\True$2021$}
		{$2027$} 
	\loigiai{Áp dụng công thức số hạng tổng quát $u_{n}=u_1+(n-1)d$ ta được  $$u_{675}=u_1+674\cdot d=-1+674\cdot3=2021.$$}
\end{ex}

% nguồn: Ok_TracNghiemLop11_Lan1
\begin{ex}%[1D2N2-4]%[ma:564fb913]
	Cho cấp số cộng $(u_n)$ có số hạng đầu $u_1=-4$ và  công sai $d=10$. Tìm số hạng thứ $37$ của cấp số cộng đã cho.
	 \choice
		{  $u_{37}=366$}
		{ $u_{37}=352$} 
		{  $u_{37}=376$}
		{ \True $u_{37}= 356$} 
	 \loigiai{
		Ta có $u_{37}=u_1+36\cdot d=  -4+36\cdot 10=356$.
	} 
\end{ex}

% nguồn: Ok_TracNghiemLop11_Lan1
\begin{ex}%[1D2N2-4]%[ma:8c965ef1]
	Cho cấp số cộng $\left(u_n \right)$ có $u_1=-0{,}1$; $d=0{,}1$. Số hạng  thứ $7$ của cấp số cộng này là
	 \choice
		 {$1{,}6$}
		 {$6$}
		 {\True $0{,}5$}
		 {$0{,}6$}
	 \loigiai{
		\[u_7=-0{,}1+6\cdot0{,}1=0{,}5.\] 
	 }
\end{ex}

% nguồn: Ok_TracNghiemLop11_Lan1
\begin{ex}%[1D2N2-4]%[ma:cdc9fb98]
	Cho cấp  số cộng có $u_1=1$, $u_2=3$. Hãy tìm công sai $d$ của cấp số cộng đó.
	\choice 
		{\True  $d=2$}
		{$d=-2$} 
		 {$d=3$}
		{$d=4$} 
	 \loigiai{
		Ta có $u_2=u_1+d\Rightarrow d=u_2-u_1=3-1=2$.  Vậy công sai $d=2$.
	 }
\end{ex}

% nguồn: Ok_TracNghiemLop11_Lan1
\begin{ex}%[1D2N2-4]%[ma:c39f7926]
	 Cho cấp số cộng có $u_1=-3$, $d=4$. Chọn khẳng định đúng trong các khẳng định sau
	 \choice
		 {$u_5=15$}
		{$u_4=8$} 
		{\True $u_3=5$} 
		{$u_2=2$} 
	 \loigiai{
		Ta  có $u_n=u_1+\left(n-1\right)d\Rightarrow u_3=u_1+2d=-3+2\cdot4=5$.
	} 
\end{ex}

% nguồn: Ok_Toan11_TracNghiem_Lan2
\begin{ex}%[1D2H2-4]%[ma:c23c4c52]
	Cho cấp số cộng $(u_n)$ với $u_1=11$ và  công sai $d=3$. Giá trị của $u_2$ bằng
	\choice 
		 {$8$}
		{$33$} 
		{$\dfrac{11}{3}$} 
		 {\True $14$}
	\loigiai{Ta có $u_2=u_1+d=11+3=14$. 
	 }
\end{ex}

% nguồn: Ok_Toan11_TracNghiem_Lan2
\begin{ex}%[1D2H2-4]%[ma:9fd33fca]
	Cho cấp số  cộng $(u_n)$ có $u_1=1$, $d=2$. Tính $u_{10}$.
	\choice 
		{$u_{10}=20$} 
		{$u_{10}=10$} 
		 {\True $u_{10}=19$}
		 {$u_{10}=15$}
	\loigiai{Ta có $u_{10}=u_1+9d=1+9  \cdot 2=19$.
	 }
\end{ex}

% nguồn: Ok_Toan11_TracNghiem_Lan2
\begin{ex}%[1D2H2-4]%[ma:c61c0b64]
	Gọi $S_n$ là tổng $n$ số hạng đầu tiên trong cấp số cộng $(a_n)$. Biết  $S_6=S_9$, tỉ số $\dfrac {a_3}{a_5}$ bằng
	 \choice
		{$\dfrac {9}{5}$} 
		{$\dfrac  {5}{9}$}
		 {\True $\dfrac {5}{3}$}
		{$\dfrac {3}{5}$} 
	\loigiai{Ta  có
		\allowdisplaybreaks 
		\begin{eqnarray*}
			 &&S_6=S_9\\
			&\Leftrightarrow& \dfrac {6\left(2a_1+5d \right)}{2}=\dfrac {9\left( 2a_1+8d \right)}{2}\\ 
			 &\Leftrightarrow& a_1=-7d.
		\end{eqnarray*}
		Tỉ  số $\dfrac {a_3}{a_5}=\dfrac {a_1+2d}{a_1+4d}=\dfrac {-7d+2d}{-7d+4d}=\dfrac {5}{3}$.
	} 
\end{ex}

% nguồn: Ok_Toan11_TracNghiem_Lan2
\begin{ex}%[1D2H2-4]%[ma:2fade1ef]
	Cho cấp số  cộng $(u_n)$, biết $u_2=3$ và $u_4=7$. Giá trị của $u_{15}$ bằng
	\choice 
		 {$27$}
		{$31$} 
		 {$35$}
		{\True $29$} 
	\loigiai{ 
		Từ giả thiết $u_2=3$ và $u_4=7$. Ta có hệ phương trình $\heva{& u_1+d=3\\ &u_1+3d=7} \Leftrightarrow \heva{&d=2\\  &u_1=1.}$\\
		 Vậy $u_{15}=u_1+14d=29$.
	} 
\end{ex}

% nguồn: Ok_Toan11_TracNghiem_Lan2
\begin{ex}%[1D2H2-4]%[ma:1ce7b519]
	Cho cấp số cộng  có $u_1=11$, $d=4$. Giá trị của số hạng $u_{99}$ là
	 \choice
		 {$401$}
		{\True  $403$}
		 {$402$}
		{$404$} 
	\loigiai{ 
		Ta có $  u_{99}=u_1+98d=11+392=403$.
	 }
\end{ex}

% nguồn: Ok_Toan11_TracNghiem_Lan2
\begin{ex}%[1D2H2-4]%[ma:a246768c]
	Cho cấp số cộng  $(u_n)$ có $u_1=1$ và $u_2=3$. Giá trị của $u_3$ bằng
	\choice 
		 {$6$}
		{$9$} 
		 {$4$}
		 {\True$5$}
	\loigiai{Ta có  $d=u_2-u_1=3-1=2$.\\
	Do đó $u_3=u_1+2d=1+2  \cdot 2 =5$.}
\end{ex}

% nguồn: Ok_Toan11_TracNghiem_Lan2
\begin{ex}%[1D2H2-4]%[ma:47629449]
	Cho cấp số cộng $(u_n)$, biết $u_1=2$ và $u_4=8$. Giá trị của  $u_5$ bằng
	\choice 
		{$12$} 
		 {\True $10$}
		 {$9$}
		 {$11$}
	\loigiai{Từ giả  thiết $u_1=2$ và $u_4=u_1+3d=8\Rightarrow d=2$.\\
		Vậy  $u_5=u_1+4d=2+4\cdot 2=10$.
	} 
\end{ex}

% nguồn: Ok_Toan11_TracNghiem_Lan2
\begin{ex}%[1D2H2-4]%[ma:6a482bdd]
	Viết ba số xen giữa $2$ và $22$ để ta được một cấp số cộng có $5$ số hạng. 
	\choice 
		{$6; 12;  18$}
		{$8;  13, 18$}
		{\True $7; 12;  17$}
		 {$6; 10; 14$}
	\loigiai{ 
		Xem cấp số cộng cần tìm là $(u_n)$ có $\heva{& u_1 = 2\\ & u_5 = 22} \Rightarrow \heva{&u_1 = 2\\& d =  5.}$\\
		Vậy cấp số cộng cần tìm  là $2; 7; 12; 17; 22$.
	} 
\end{ex}

% nguồn: Ok_Toan11_TracNghiem_Lan2
\begin{ex}%[1D2H2-4]%[ma:f04b204c]
	Một cấp số cộng $(u_n)$ có $u_{13}=8$ và $d=-3$. Số hạng thứ mười ba của cấp  số cộng $(u_n)$ là
	 \choice
		 {$50$}
		 {$28$}
		{\True  $38$}
		{$44$} 
	\loigiai{ 
		Ta  có
		\allowdisplaybreaks 
		\begin{eqnarray*}
			&& u_{13}=u_1+12d  \Leftrightarrow 8=u_1+12\cdot(-3)\\
			 & \Leftrightarrow u_1&=44 \Leftrightarrow {u_3}=u_1+2d=44-6=38.
		\end{eqnarray*}
	 }
\end{ex}

% nguồn: Ok_Toan11_TracNghiem_Lan2
\begin{ex}%[1D2H2-4]%[ma:b60f08f0]
	Cho  dãy số $(u_n)$ thỏa mãn $u_9=5u_2$ và $ u_{13}=2 u_6+5$.Giá trị của số hạng đầu $u_1$ và công sai $d$ của cấp số cộng là
	 \choice
		{\True $u_1=3,  \ d=4$}
		 {$u_1=3, \ d=5$}
		 {$u_1=4, \ d=5$}
		{$u_1=4, \  d=3$}
	 \loigiai{
		Sử dụng công thức  $u_n=u_1+(n-1) d$.\\ Ta có
		\[ 
		 \heva{&
		u_1+8d=5(u_1+d)  \\
		 &u_1+12d=2(u_1+5d)+5.}
		\] 
		Giải  hệ:
		\[ 
		\heva{& 
		 4u_1-3d=0 \\&
		u_1-2d=-5.} 
		\Rightarrow u_1=3, \  d=4.
		\] 
	} 
\end{ex}

% nguồn: Ok_Toan11_TracNghiem_Lan2
\begin{ex}%[1D2H2-4]%[ma:441d4d5f]
	Cho cấp số cộng $(u_n)$ với $u_1=2$; $d=9$. Khi đó  số $2\,018$ là số hạng thứ mấy trong dãy.
	\choice 
		{$226$} 
		{\True  $225$}
		{$223$} 
		{$224$} 
	\loigiai{ 
		$u_n=u_1+(n-1)\cdot d \Leftrightarrow 2\,018=2+(n-1)\cdot 9 \Leftrightarrow  n=225$.
	} 
\end{ex}

% nguồn: Ok_Toan11_TracNghiem_Lan2
\begin{ex}%[1D2H2-4]%[ma:4148994d]
	Cho cấp số cộng $1;4;7;\ldots$ Số hạng thứ $100$  của cấp số cộng là
	 \choice
		{$297$} 
		 {$301$}
		{$295$} 
		{\True  $298$}
	 \loigiai{
		Cấp số cộng $1;4;7;\ldots$ có số hạng đầu  $u_1=1$ và công sai $d=3$.\\
		Số hạng thứ  $100$ của cấp số cộng là $u_{100}=u_1+99\cdot d=1+99\cdot 3=298$.
	} 
\end{ex}

% nguồn: Ok_Toan11_TracNghiem_Lan2
\begin{ex}%[1D2H2-4]%[ma:786e9520]
	 Cho cấp số cộng $(u_n)$ biết $u_1=3$, $u_8=24$ thì $u_{11}$ bằng
	 \choice
		{$30$} 
		 {\True $33$}
		{$32$} 
		{$28$} 
	\loigiai{ 
		Ta  có
		 $u_8=u_1+7d \Rightarrow d=\dfrac{u_8-u_1}{7}=\dfrac{24-3}{7}=3$.\\
		Khi đó  $u_{11}=u_1+10d=33$.
	 }
\end{ex}

% nguồn: Ok_Toan11_TracNghiem_Lan2
\begin{ex}%[1D2H2-4]%[ma:16bc9e52]
	Cho cấp số cộng  có
	 $u_1=-3$
	, 
	$d=4$ 
	. Khẳng  định nào sau đây đúng?
	 \choice
		{${u_5}=15$} 
		 {${u_4}=8$}
		{\True  ${u_3}=5$}
		 {$u_2=2$}
	 \loigiai{
		Ta có 
		$u_3=u_1+2d=-3+2\cdot  4=5$.
	 }
\end{ex}

% nguồn: Ok_Toan11_TracNghiem_Lan2
\begin{ex}%[1D2H2-4]%[ma:9c5e63ae]
	Cấp  số cộng $(u_n)$ có số hạng đầu $u_1=3$, công sai $d=-2$ thì số hạng thứ $5$ là
	 \choice
		{$u_5=1$} 
		{$u_5=8$} 
		{$u_5=-7$} 
		{\True $u_5=-5$} 
	\loigiai{Ta có $u_5=u_1+4d=3+4 \cdot (-2)=-5$. 
	} 
\end{ex}

% nguồn: Ok_Toan11_TracNghiem_Lan2
\begin{ex}%[1D2H2-4]%[ma:89ff2f01]
	Cho  cấp số cộng có $u_1=3$, $d=-2$. Giá trị của $u_5$ là
	\choice 
		{$u_5=8$} 
		 {$u_5=1$}
		 {\True $u_5=-5$}
		 {$u_5=-7$}
	 \loigiai{
		Ta  có $
		u_5=u_1+4d=3+4\cdot(-2)=-5$. 
	} 
\end{ex}

% nguồn: Ok_Toan11_TracNghiem_Lan2
\begin{ex}%[1D2H2-4]%[ma:e139d335]
	Giải phương trình $1+8+15+22+\ldots +x=7\,944$. 
	\choice 
		{\True $x=330$} 
		 {$x=220$}
		{$x=351$} 
		{$x=407$} 
	\loigiai{Ta có cấp số cộng  với $u_1=1$, $d=7$, $u_n=x$, $S_n=7\,944$.\\
		Áp dụng  công thức
		\allowdisplaybreaks 
		\begin{eqnarray*}
			&&S_n=\dfrac {\left[2u_1+\left(n-1\right)d \right]n}{2}\\ 
			 &\Leftrightarrow&7\,944=\dfrac {\left[2\cdot 1+\left(n-1 \right)7\right]n}{2}\\
			 &\Leftrightarrow&7n^2-5n-15\,888=0\\
			&\Leftrightarrow&\hoac{&n=48\ \left(\text{thỏa} \right) \\  & n=-\dfrac {331}{7}\ \left(\text{loại}\right).}
		\end{eqnarray*}
		Vậy $x=u_{48}=1+47\cdot 7=330$. 
	} 
\end{ex}

% nguồn: Ok_Toan11_TracNghiem_Lan2
\begin{ex}%[1D2H2-4]%[ma:7d5224e3]
	Cho cấp số cộng có $u_1=-2$ và $d=4$.  Chọn khẳng định đúng trong các khẳng định sau.
	 \choice
		 {$u_4=8$}
		{$u_5=15$} 
		 {$u_2=3$}
		 {\True $u_3=6$}
	\loigiai{ 
		 Ta có $u_1=-2$ và $d=4$. Suy ra\\
		$u_2=u_1+d=-2+4=2$;  $u_3=u_1+2\cdot d=-2+2\cdot 4=6$; $u_4=u_1+3\cdot d=-2+3\cdot 4=10$; $u_5=u_1+4\cdot d=-2+4\cdot 4=14$.
	 }
\end{ex}

% nguồn: Ok_Toan11_TracNghiem_Lan2
\begin{ex}%[1D2H2-4]%[ma:d82e1ebf]
	Cho cấp số  cộng $(u_n)$ có $u_3=2$, công sai $d=-2$. Số hạng thứ hai của cấp số cộng đó là
	\choice 
		 {\True $u_2=4$}
		 {$u_2=0$}
		{$u_2=-4$} 
		{$u_2=3$} 
	\loigiai{Ta có $d=u_3-u_2 \Rightarrow -2=2-u_2 \Rightarrow u_2=4$. 
	} 
\end{ex}

% nguồn: Ok_Toan11_TracNghiem_Lan2
\begin{ex}%[1D2H2-4]%[ma:6d167826]
	Cho cấp số cộng $(u_n)$  biết $u_1=3$, $u_2=-1$. Giá trị của số hạng $u_3$ là
	\choice 
		{$u_3=4$} 
		{$u_3=7$} 
		{$u_3=2$} 
		{\True $u_3=-5$} 
	\loigiai{ 
		Dãy cấp số cộng thỏa  mãn $2u_2=u_1+u_3$ nên
		\[ 
		u_3=2u_2-u_1=2\cdot(-1)-3=-5. 
		\] 
	 }
\end{ex}

% nguồn: Ok_Toan11_TracNghiem_Lan2
\begin{ex}%[1D2H2-4]%[ma:14860004]
	Cho cấp số cộng $(u_n)$ có $u_2=2001$  và $u_5=1995$. Khi đó $u_{1001}$ bằng
	\choice 
		{$4005$} 
		{$1$} 
		{\True  $3$}
		 {$4003$}
	\loigiai{ 
		Gọi $u_1$ và  $d$ lần lượt là số hạng đầu tiên và công sai của cấp số cộng.\\
		Ta có $\heva{& u_2=2001\\ &u_5=1995} \Leftrightarrow \heva{& u_1+d=2001\\ &u_1+4d=1995} \Leftrightarrow  \heva{& u_1=2003\\ &d=-2.}$\\
		Vậy  $u_{1001}=u_1+1000d=3$.
	} 
\end{ex}

% nguồn: Ok_Toan11_TracNghiem_Lan2
\begin{ex}%[1D2H2-4]%[ma:c26d656d]
	Cho cấp số cộng $(u_n)$ có $u_1=123$ và $u_8-u_{15}=84$. Số $11$ là số hạng  thứ bao nhiêu của cấp số cộng đã cho?
	 \choice
		 {\True $17$}
		 {$16$}
		 {$18$}
		 {$19$}
	\loigiai{ 
		 Ta có $u_8-u_{15}=84\Leftrightarrow u_1+2d-(u_1+14d)=84\Leftrightarrow d=-7$.\\
		Số hạng tổng  quát $u_n=-7n+130$.\\
		Ta có $u_n=11\Leftrightarrow n=17$. 
	} 
\end{ex}

% nguồn: Ok_TracNghiemLop11_Lan1
\begin{ex}%[1D2H2-4]%[ma:d7fec897]
	Cho cấp số  cộng $(u_n)$ có số hạng đầu $u_1 = -2$ và công sai $d = 3$. Số $-168$ là số hạng thứ bao nhiêu?
	 \choice
		 {$28$}
		 {$30$}
		{\True $29$} 
		 {$31$}
	 \loigiai{
		Ta có \\ 
		\begin{align*}
			& u_n = u_1(n-1)d  \\
			\Leftrightarrow  & -168=-2\cdot (n-1)\cdot 3.\\
			\Leftrightarrow & n = \dfrac{-168}{-6} + 1  = 29.
		\end{align*}
	} 
\end{ex}

% nguồn: Ok_TracNghiemLop11_Lan1
\begin{ex}%[1D2H2-4]%[ma:2cd6ea8c]
	Cấp số cộng $(u_n)$ có  $u_6=12$, $u_{10}=24$. Tìm số hạng đầu $u_1$.
	\choice 
		{$u_1=3$} 
		{$u_1=2$} 
		{$u_1=5$} 
		{\True $u_1=-3$} 
	\loigiai{Giả sử cấp số cộng có công sai là $d$. Ta  có hệ
		\begin{center}
			$\heva{&u_6=u_1+5d\\&u_{10}=u_1+9d}\Leftrightarrow\heva{&12=u_1+5d\\&24=u_1+9d}\Leftrightarrow\heva{&u_1=-3\\&d=3.}$ 
		\end{center}
	} 
\end{ex}

% nguồn: Ok_TracNghiemLop11_Lan1
\begin{ex}%[1D2H2-4]%[ma:b4c892f1]
	Cho cấp số cộng $\left(u_n \right)$ có số hạng đầu $u_1=-5$ và công sai $d=3$. Số $100$ là số hạng thứ mấy của cấp  số cộng?
	 \choice
		{$15$} 
		 {$20$}
		 {$35$}
		 {\True $36$}
	\loigiai{ 
		Ta có $u_n=u_1+\left(n-1\right)d\Leftrightarrow 100=-5+\left(n-1\right)\cdot3\Leftrightarrow 100=3n-8\Leftrightarrow n=36$. 
	} 
\end{ex}

% nguồn: Ok_TracNghiemLop11_Lan1
\begin{ex}%[1D2H2-4]%[ma:1e83dd05]
	Cho cấp số cộng $\left(u_n \right)$ có $u_1=1$,  $d=4$. Tìm số hạng $u_{12}$.
	\choice 
		{$u_{12}=31$} 
		 {$u_{12}=13$}
		{\True  $u_{12}=45$}
		 {$u_{12}=17$}
	\loigiai{ 
		Ta có  $u_{12}=u_1+11d=45$.
	} 
\end{ex}

% nguồn: Ok_TracNghiemLop11_Lan1
\begin{ex}%[1D2H2-4]%[ma:f5afd1ce]
	Cho hai số $-3$ và $23$. Xen kẽ giữa hai số đã cho $n$ số hạng để tất cả các số đó tạo thành cấp số  cộng có công sai $d=2$. Tìm $n$.
	 \choice
		{\True  $n=12$}
		 {$n=13$}
		 {$n=14$}
		 {$n=15$}
	 \loigiai{
		Theo giả thiết thì ta được một cấp số cộng có $n+2$ số hạng với $u_1=-3$,  $u_{n+2}=23$.\\
		 Khi đó $u_{n+2}=u_1+\left(n+1\right)d\Leftrightarrow n+1=\dfrac{u_{n+2}-u_1}{d}=\dfrac{23-\left(-3\right)}{2}=13\Leftrightarrow n=12$.
	 }
\end{ex}

% nguồn: Ok_TracNghiemLop11_Lan1
\begin{ex}%[1D2H2-4]%[ma:11cdc63d]
	Cho cấp số  cộng $\left(u_n \right)$ có $u_1=2$, $d=3$. Số hạng thứ $7$ của cấp số cộng này là
	\choice 
		{\True $20$} 
		 {$21$}
		{$19$} 
		{$23$} 
	 \loigiai{
		Ta có số  hạng tổng quát của cấp số cộng là $u_n=u_1+\left(n-1\right)d\Rightarrow u_7=2+6\cdot3=20$.
	} 
\end{ex}

% nguồn: Ok_TracNghiemLop11_Lan1
\begin{ex}%[1D2H2-4]%[ma:b17fdc72]
	Cho cấp số  cộng $(u_n)$ có $u_1=4$; $u_2=1$. Giá trị của $u_{10}$ bằng
	 \choice
		 {$u_{10}=-31$}
		{\True $u_{10}=-23$} 
		{$u_{10}=-20$} 
		 {$u_{10}=15$}
	 \loigiai{
		 Ta có $u_2=u_1+d\Rightarrow d=-3$.\\
		Khi đó  $u_{10}=u_1+9d\Leftrightarrow u_{10}=4+9\cdot(-3)\Leftrightarrow u_{10}=-23$.
	 }
\end{ex}

% nguồn: Ok_TracNghiemLop11_Lan1
\begin{ex}%[1D2C2-4]%[ma:de7ec5ad]
	\immini[thm]{Trong hội chợ, một công ty  sơn muốn xếp $169$ hộp sơn theo số lượng $1$, $ 3$, $ 5$, $ \ldots$ từ trên xuống dưới (như hình vẽ). Hỏi hàng cuối cùng có bao nhiêu hộp sơn?
	\choice[2] 
		{$13$} 
		 {\True $25$}
		 {$27$}
		{$20$}} 
		{ 
	\definecolor{amber}{rgb}{1.0, 0.75, 0.0} 
	\definecolor{ceruleanblue}{rgb}{0.16, 0.32,  0.75}
	 \definecolor{alizarin}{rgb}{0.82, 0.1, 0.26}
	\definecolor{deepcarminepink}{rgb}{0.94, 0.19, 0.22} 
	\definecolor{melon}{rgb}{0.99, 0.74,  0.71}
	\definecolor{azure(colorwheel)}{rgb}{0.0, 0.5, 1.0} 
	\definecolor{gainsboro}{rgb}{0.86,  0.86, 0.86}
		\begin{tikzpicture}[>=stealth,x=1cm,y=1cm,scale=0.6]
				\draw[->,black]  (-2.2,0) -- (4,0) node[below] {$x$};
				\foreach \x in{-2,3} \draw[xshift=\x cm] (0pt,-2pt)--(0pt,2pt); 
				\draw[->,black] (0,-2.5) --  (0,3.5) node[left] {$y$};
				\foreach \y in{-2,-1,2,3} \draw[yshift=\y cm] (-2pt,0pt)--(2pt,0pt); 
				 \draw [-,black](-1,-2.5)--(-1,3.5);
				\draw  (1,0) node[below]{$1$} circle (1pt);
				\draw (2,0) node[below]{$2$}  circle (1pt);
				\draw (-1,0)  node[above left]{$-1$} circle (1pt);
				\draw  (0,1) node[left]{$1$} circle (1pt);
				\draw (0,0) node[below  right]{$O$} circle (1pt);
				\draw [dashed]  (1,0)--(1,1)--(0,1);
				\draw[smooth,samples=100,domain=-0.8:3] plot  (\x,{ln((\x)+1)/ln(2)});
			\end{tikzpicture}1
		} 
	\loigiai{ 
		 Số hộp sơn được xếp theo một cấp số cộng với công sai $d=2$ và $u_1=1$.\\
		Với tổng của cấp số cộng là $S_n=\dfrac{n\left[2u_1+(n-1)d\right]}{2}=169 \Rightarrow  n=13$.\\
		Vậy, hàng cuối cùng là hàng thứ $13$ với  số hộp sữa được xếp
		$u_{13}=u_1+(n-1)\cdot d=1+(13-1)\cdot  2=25$.
	} 
\end{ex}


%%%%%%%%%%%%%%%% ĐÚNG/SAI %%%%%%%%%%%%%%%%

% nguồn: Ok_dungsai_lop11_Lan1
\begin{ex}%[1D2N2-4]%[ma:caf02041]
 Cho dãy số $(u_n)=\heva{&u_1=3\\&u_n =u_{n-1}+5},\,\forall n\in\mathbb{N},n\ge  2$. Các mệnh đề sau đúng hay sai?
  \choiceTF
  {\True Dãy số có số hạng  thứ $2$ là $u_2=8$}
  {\True  Dãy số $(u_n)$ là một cấp số cộng}
  {\True Dãy số $(u_n)$ là  dãy tăng}
  {Dãy số $(u_n)$ bị chặn trên} 
  \loigiai{
  \begin{itemchoice}
   \itemch \textbf{Đúng.} Ta  có $u_2=u_1+5=3+5=8$.
   \itemch \textbf{Đúng.} Ta có $(u_n)$ là cấp số cộng với công  sai $d=5$.
   \itemch \textbf{Đúng.} Ta có $u_n-u_{n-1}=5>0\,(\forall n\in\mathbb{N},n\ge 2)\Leftrightarrow u_n  >u_{n-1}\,(\forall n\in\mathbb{N},n\ge 2)$ nên $(u_n)$ là dãy tăng.
   \itemch \textbf{Sai.} Ta có $(u_n)$ là cấp số cộng với công sai $d=5>0$ nên dãy không bị chặn  trên.
  \end{itemchoice}
 } 
\end{ex}

% nguồn: Ok_dungsai_lop11_Lan2
\begin{ex}%[1D2N2-4]%[1D2N2-6]%[ma:066bf7bd]
	Cho cấp số cộng $(u_n)$ có số hạng đầu $u_1=\dfrac{3}{2}$, công sai  $d=\dfrac{1}{2}.$ Khi đó
	 \choiceTF
		{Công thức  cho số hạng tổng quát $u_n=1+\dfrac{n}{3}$}
		{\True $5$ là số hạng thứ  $8$ của cấp số cộng đã cho}
		{$\dfrac{15}{4}$ là một số hạng của cấp số cộng  đã cho}
		{Tổng $100$ số hạng đầu của  cấp số cộng $(u_n)$ bằng $2\,620$}
	\loigiai{ 
		\begin{itemchoice}
			\itemch {\bf Sai}.\\ Ta có $u_n=u_1+(n-1)  d=\dfrac{3}{2}+(n-1) \cdot \dfrac{1}{2}=1+\dfrac{n}{2}$.
			\itemch {\bf Đúng}.\\ Xét $5=1+\dfrac{n}{2} \Rightarrow n=8\in  \mathbb{N}^*$. Suy ra $5$ là số hạng thứ $8$ của cấp số cộng đã cho.
			\itemch {\bf Sai}.\\ Xét $\dfrac{15}{4}=1+\dfrac{n}{2} \Rightarrow n=\dfrac{11}{2} \notin \mathbb{N}^*$. Suy ra $\dfrac{15}{4}$ không phải là  một số hạng của cấp số cộng đã cho.
			\itemch {\bf Sai}.\\ Tổng $100$ số hạng  đầu của cấp số cộng là\\
			$S_{100}=\dfrac{100\left[2\cdot  \dfrac{3}{2}+(100-1) \cdot \dfrac{1}{2}\right]}{2}=2\,625$.
		\end{itemchoice}
	 }
\end{ex}

% nguồn: Ok_dungsai_lop11_Lan2
\begin{ex}%[1D2N2-4]%[1D2N2-6]%[ma:4d08dc73]
	Cho cấp số cộng $(u_n)$ có  $u_1=-5$, công sai $d=3$.
	 \choiceTF
		{\True Số $100$ là số hạng thứ 36 của  cấp số cộng}
		{Số hạng thứ 3 của cấp  số cộng bằng $5$}
		{Tổng 10 số hạng đầu tiên của cấp số cộng  bằng $250$}
		{\True Kể từ số hạng thứ 3 thì các số hạng của cấp số cộng đều nhận  giá trị dương}
	 \loigiai{
		Ta có 
		\begin{itemize}
			\item $100=u_1+(n-1)d=-5+(n-1)\cdot  3\Leftrightarrow n=36$.
			 \item $u_3=u_1+2d=-5+2\cdot 3=1$.
			\item  $S_{10}=\dfrac{\left(u_1+u_{10}\right)10}{2}=\dfrac{\left(u_1+u_1+9d\right)\cdot 10}{2}=5(-10+9\cdot 3)=85$.
			\item $u_n=u_1+(n-1)d=-5+(n-1)\cdot 3=3n-8> 0\Leftrightarrow n  > \dfrac{8}{3}$.
		\end{itemize}
		\begin{itemchoice}
			 \itemch {\bf Đúng}.\\
			\itemch {\bf  Sai}.\\
			\itemch {\bf Sai}.\\ 
			\itemch  {\bf Đúng}.\\
		\end{itemchoice}	.
	} 
\end{ex}

% nguồn: Ok_dungsai_lop11_Lan2
\begin{ex}%[1D2H2-4]%[ma:29c8f7aa]
	Cho  cấp số cộng $(u_n)$ có $u_1 = -5$, công sai $d = 3$. Xét tính đúng sai của các khẳng định sau
	\choiceTF 
		{\True Số $100$ là số hạng thứ $36$ của  cấp số cộng}
		{Số hạng thứ $3$  của cấp số cộng bằng $5$}
		{Tổng $10$ số hạng đầu tiên của cấp số cộng bằng $250$} 
		{\True Kể từ số hạng thứ $3$ thì các số hạng của cấp số cộng đều nhận giá trị  dương}
	 \loigiai{
		Ta có $100 = u_1 +  (n-1)d = -5 + (n-1)\cdot 3 \Leftrightarrow n=36$.\\
		Ta có 
		$u_3 = u_1 + 2d = -5 + 2\cdot  3 = 1$.\\
		Suy  ra
		\begin{eqnarray*}
			S_{10} &=& \dfrac{(u_1 +  u_{10})10}{2}\\
			& =& \dfrac{(u_1 + u_1 + 9d)\cdot 10}{2}\\ 
			& =&5(-10+9\cdot  3) = 85.
		\end{eqnarray*}
		 Và
		\[u_n = u_1 + (n-1)d = -5+(n-1)\cdot  3 = 3n-8 > 0 \Leftrightarrow n > \dfrac{8}{3}.\]
		\begin{itemchoice}
			\itemch Số $100$ là số hạng thứ $36$ của  cấp số cộng.
			\itemch Số hạng thứ $3$ của  cấp số cộng bằng $1$.
			\itemch Tổng $10$ số hạng đầu tiên của cấp số cộng bằng $85$. 
			\itemch Kể từ số hạng thứ $3$ thì các số hạng của cấp số cộng đều nhận giá trị dương. 
		\end{itemchoice}
	 }
\end{ex}

% nguồn: Ok_dungsai_lop11_Lan2
\begin{ex}%[1D2H2-4]%[ma:3998d5de]
	Cho cấp số cộng $(u_n)$ có $u_1=5$ và $d=-7$. Xét tính đúng sai của các  khẳng định sau.
	\choiceTF 
		{\True $u_{11}=-65$} 
		 {$u_5+u_7=-50$}
		{\True Số $-849$ là số hạng thứ $123$ của  cấp số cộng}
		{\True Số $-114$ là số hạng thứ $18$ của cấp số  cộng}
	\loigiai{ 
		Số hạng tổng quát của cấp số cộng là  $u_n=u_1+(n-1)\cdot d=5+(n-1)\cdot (-7)=-7n+12$
		\begin{itemchoice}
			 \itemch
			Vì ta có $u_{11}=-7\cdot 11+12=-65$. 
			\itemch 
			Vì $u_5+u_7=(u_1+4\cdot d)+(u_1+6\cdot d)=2u_1+10\cdot  d=2\cdot 5+10\cdot (-7)=10-70=-60$.
			\itemch 
			Vì Ta  có $-849=-7n+12 \Rightarrow -7n=-861 \Rightarrow n=123$.
			 \itemch
			Vì Ta có $-114=-7n+12 \Rightarrow -7n=-126  \Rightarrow n=18$.
		\end{itemchoice}
	 }
\end{ex}

% nguồn: Ok_dungsai_lop11_Lan2
\begin{ex}%[1D2H2-4]%[ma:e3e9915a]
	Cho cấp số cộng $(u_n)$ có $u_1=123$, $u_3 - u_{15} =  84$. Xét tính đúng sai của các khẳng định sau
	\choiceTF 
		{\True Số hạng thứ  $17$ của cấp số cộng là $u_{17}=11$}
		{\True Công  sai của cấp số cộng là $d=-7$}
		 {Số hạng thứ $2$ của cấp số cộng là $u_2=130$}
		 {Tổng $17$ số hạng đầu tiên là $S_{17}=1\,130$}
	\loigiai{ 
		Giả sử  cấp số cộng $(u_n)$ có công sai $d$.\\
		Theo giả thiết ta có $u_3 - u_{15} = 84 \Leftrightarrow u_1 + 2d - u_1 - 14d  = 84 \Leftrightarrow -12d = 84 \Leftrightarrow d = -7$.\\
		Vậy $u_{17} = u_1 + 16d = 123 + 16(-7) = 11$; $u_2 = u_1 + d = 123 +  (-7) = 116$.\\
		\[S_{17} = \dfrac{(u_1 + u_{17})17}{2} = \dfrac{(123  + 11)17}{2} = 1\,139.\]
		\begin{itemchoice}
			\itemch Số hạng thứ $17$ của cấp số cộng  là $u_{17}=11$.
			\itemch Công sai  của cấp số cộng là $d=-7$.
			\itemch Số hạng  thứ $2$ của cấp số cộng là $u_2=116$.
			\itemch Tổng $17$ số hạng  đầu tiên của cấp số cộng là $S_{17}=1\,139$.
		\end{itemchoice}
	 }
\end{ex}

% nguồn: Ok_dungsai_lop11_Lan2
\begin{ex}%[1D2H2-4]%[ma:207b7ab6]
	
	Cho  cấp số cộng $\left(u_n\right)$ có công sai $d < 0$ thoả mãn ${\heva{&u_1+u_7=26 \\ &u_2^2+u_6^2=466}}$. Khi đó:
	 \choiceTF
		 {\True Số hạng $u_1=25$}
		{Công sai $d=-3$} 
		 {\True Số hạng $u_{10}=-11$}
		{\True  Số hạng $u_{2024}=-8067$}
	 \loigiai{
		Ta có $\heva{& {u_1}+{u_7}=26 \\ & u_2^2+u_6^2=466 }$ $\Leftrightarrow $ $\heva{& 2{u_1}+6d=26 \\ & {{( {u_1}+d)}^2}+{{( {u_1}+5d)}^2}=466 }$ $\Leftrightarrow $ $\heva{& {u_1}=13-3d( 1) \\ & {{( {u_1}+d)}^2}+{{( {u_1}+5d)}^2}=466( 2)  }$.\\
		Thay $( 1)$ vào $( 2)$, ta được ${(13-2  d)^2+(13+2 d)^2=466 \Leftrightarrow 8 d^2+338=466}$\\
		$\Leftrightarrow $ $\hoac{&  d=4 \\ & d=-4 }$\\
		Vì $ d<0$ nên ta nhận $ d=-4$,  khi đó ${u_1}=25$.\\
		Khi đó  ${u_n=u_1+(n-1) d=25+(n-1)(-4)=29-4 n}$.
		\begin{itemchoice}
			 \itemch \textbf{Đúng.}
			 \itemch \textbf{Sai.}
			\itemch  \textbf{Đúng.}\\
			$u_{10}=29-4\cdot10=-11$. 
			\itemch \textbf{Đúng.}\\ 
			$u_{2024}=29-4\cdot2024=-8067$. 
		\end{itemchoice}
	 }
\end{ex}

% nguồn: Ok_dungsai_lop11_Lan2
\begin{ex}%[1D2H2-4]%[ma:c127250a]
	Một cấp  số cộng có năm số hạng mà tổng số hạng đầu và số hạng thứ tư bằng $36$, tổng của số hạng thứ hai và số hạng cuối bằng $44$.
	\choiceTF 
		{Dãy cấp số cộng có số hạng đầu  $u_1=11$}
		{\True Dãy cấp số cộng  có tổng $u_1+u_5=40$}
		{\True Dãy cấp số cộng có  $u_2=16$}
		{Tổng của $3$ số hạng đầu  tiên trong dãy cấp số cộng bằng $45$}
	\loigiai{ 
		 Ta có
		\begin{eqnarray*}
			& & \heva{&u_1+u_4=36\\&u_2+u_5=44}\\ 
			&\Leftrightarrow  & \heva{&u_1+\left(u_1+3d\right)=36\\&\left(u_1+d\right)+\left(u_1+4d\right)=44}\\
			&\Leftrightarrow  & \heva{&2u_1+3d=36\\&2u_1+5d=44}\\
			&\Leftrightarrow  & \heva{&u_1=12\\&d=4.}
		\end{eqnarray*}
		 Suy ra $\heva{&u_2=u_1+d=16\\&u_1+u_5=u_1+u_1+4d=40\\&S_3=\dfrac{3\left(u_1+u_3\right)}{2}=48.}$
		\begin{itemchoice}
			\itemch {\bf  Sai}.\\
			 \itemch {\bf Đúng}.\\
			\itemch  {\bf Đúng}.\\
			\itemch  {\bf Sai}.\\
		\end{itemchoice}
	} 
\end{ex}

% nguồn: Ok_dungsai_lop11_Lan2
\begin{ex}%[1D2H2-4]%[ma:ec52b8fb]
	Cho bốn số $-2$, $x$, $6$, $y$ theo  thứ tự lập thành một cấp số cộng. Khi đó
	 \choiceTF
		{\True  $x=2$}
		 {$y=8$}
		{$P=y-x=6$} 
		{\True $Q=x^2+y^2=104$} 
	 \loigiai{
		Theo tính chất của  cấp số cộng, ta có $x=\dfrac{-2+6}{2}=2$ và $6=\dfrac{x+y}{2}$.\\
		 $\Rightarrow$ $6=\dfrac{2+y}{2}\Rightarrow y=10$.\\
		\begin{itemchoice}
			 \itemch \textbf{Đúng.}
			\itemch \textbf{Sai.} 
			\itemch  \textbf{Sai.}\\
			 $P=y-x=8$.
			\itemch  \textbf{Đúng.}\\
			 $Q=x^2+y^2=2^2+10^2=104$.
		\end{itemchoice}
	} 
\end{ex}

% nguồn: Ok_dungsai_lop11_Lan2
\begin{ex}%[1D2H2-4]%[ma:58d6d696]
	Cho cấp số cộng $(u_n)$ có công sai $d < 0$ thoả mãn $\heva{& u_1+u_7=26\\& u_2^2+u_6^2=466}$. Xét tính đúng sai của các khẳng định sau 
	\choiceTF 
		{\True Số hạng $u_1=25$} 
		{Công  sai $d=-3$}
		{\True Số  hạng $u_{10}=-11$}
		{\True  Số hạng $u_{2024}=-8067$}
	 \loigiai{
		\begin{itemchoice}
			\itemch  Ta có
			\begin{align*}
				&\heva{& u_1+u_7  = 26\\& u_2^2+u_6^2 = 466}\\
				 \Leftrightarrow &\heva{& 2u_1+6\cdot d = 26\\&(u_1+d)^2+(u_1+5\cdot d)^2 = 466}\\
				\Leftrightarrow &\heva{& u_1 = 13-3d \quad (1)\\& (u_1+d)^2+(u_1+5\cdot d)^2 =  466. \quad (2)}
			\end{align*}
			Thay $(1)$ vào $(2)$  ta được
			$$(13-2\cdot d)^2+(13+2\cdot d)^2 = 466 \Leftrightarrow 8d^2+338 = 466 \Leftrightarrow \heva{& d  = 4\\ & d = -4.}$$
			Vì  $d < 0$ nên ta nhận $d = -4$, khi đó $u_{1} = 25$.
			\itemch  Vì $d=-4$
			\itemch Vì  $u_n = u_1+(n-1)\cdot d = 25+(n-1)\cdot (-4) = 29-4n$.\\
			Nên  số hạng $u_{10} = -11$.
			\itemch Số hạng $u_{2024} = -8\,067$. 
		\end{itemchoice}
	 }
\end{ex}

% nguồn: Ok_dungsai_lop11_Lan2
\begin{ex}%[1D2H2-4]%[ma:94966c90]
	Cho cấp số cộng $(u_n)$ có $u_4=-12$; $u_{14}=18$. Xét tính  đúng sai của các khẳng định sau
	 \choiceTF
		 {\True Công sai của cấp số cộng là $d=3$}
		{Số hạng đầu của cấp số cộng là  $u_1=21$}
		{\True Số hạng thứ $9$ của cấp số cộng là  $u_9=3$}
		{Tổng $5$ số hạng đầu của cấp số cộng là $S_5=-60$} 
	\loigiai{ 
		Gọi $u_1$ là số hạng đầu, $d$ là công  sai của cấp số cộng $(u_n)$.\\
		 Ta có
		\begin{eqnarray*}
			&&\heva{ &u_4 = u_1 + 3d\\ &u_{14} = u_1 + 13d}\\ &\Leftrightarrow& \heva{& u_1 + 3d = -12\\& u_1 + 13d = 18} \\ 
			 & \Leftrightarrow& \heva{&d = 3\\& u_1 = -21.}
		\end{eqnarray*}
		Số hạng thứ $9$ là $u_9  = u_1 + 8d = -21 + 8\cdot 3 = 3$.\\
		Tổng $5$ số hạng  đầu là $S_5 = 5u_1 + \dfrac{5(5-1)}{2}d = -75$.
		\begin{itemchoice}
			\itemch Công  sai của cấp số cộng là $d=3$.
			\itemch Số hạng đầu của cấp số cộng là $u_1 = -21\ne 21$. 
			 \itemch Số hạng thứ $9$ của cấp số cộng là $u_9=3$.
			\itemch Tổng $5$ số hạng đầu của cấp số  cộng là $S_5=-75\ne -60$.
		\end{itemchoice}
	 }
\end{ex}

% nguồn: Ok_dungsai_lop11_Lan2
\begin{ex}%[1D2H2-4]%[ma:b18beeaf]
	Cho cấp số cộng $\left(u_n\right)$ , biết rằng: $u_1=5$ và tổng của $50$ số hạng  đầu bằng $5150$, khi đó:
	\choiceTF 
		{Công sai của cấp số cộng bằng  $6$}
		{\True  Số hạng $u_{85}=341$}
		{Số hạng  $u_{10}=42$}
		{\True Tổng của $85$ số hạng  đầu $S_{85}=14705$}
	 \loigiai{
		Ta có $S_{50}=\dfrac{50}{2}\left(2u_1+49d\right)=\dfrac{50}{2}(2\cdot 5+49d)=5150\Rightarrow d=4$.\\ 
		Suy ra  $u_n=u_1+(n-1)d=5+(n-1)4=1+4n$.\\
		 Khi đó:
		\begin{itemchoice}
			 \itemch \textbf{Sai.}
			\itemch  \textbf{Đúng.}\\
			$u_{85}=1+4\cdot85=341$. 
			 \itemch \textbf{Sai.}\\
			 $u_{10}=1+4\cdot10=41$.
			\itemch  \textbf{Đúng.}\\
			$S_{85}=\dfrac{85}{2}\left(2u_1+84d\right)=\dfrac{85}{2}(2\cdot  5+84\cdot 4)=14705$.
		\end{itemchoice}
	 }
\end{ex}

% nguồn: Ok_dungsai_lop11_Lan2
\begin{ex}%[1D2H2-4]%[ma:b7267277]
	
	Cho cấp  số cộng $\left(u_n\right)$ thoả mãn ${\heva{&u_1-u_3+u_5=15 \\ &u_1+u_6=27}}$. Khi đó
	\choiceTF 
		{\True Số  hạng $u_1=21$}
		{Công  sai của cấp số cộng bằng $-2$}
		{\True Số hạng $u_{11}=-9$} 
		 {\True Số $-6048$ là số hạng thứ $2024$}
	\loigiai{ 
		Áp dụng công thức số hạng tổng quát của cấp số cộng ${u_n=u_1+(n-1) d}$.\\ 
		Khi đó $\heva{& {u_1}-{u_3}+{u_5}=15 \\ & {u_1}+{u_6}=27 }$ $\Leftrightarrow $ $\heva{& {u_1}-( {u_1}+2d)+( {u_1}+4d)=15 \\ & {u_1}+( {u_1}+4d)=27 }$ $\Leftrightarrow $ $\heva{& {u_1}+2d=15 \\ & 2{u_1}+5d=27 }$ $\Leftrightarrow $ $\heva{& {u_1}=21 \\ &  d=-3 }$.\\
		Suy  ra ${u_n}={u_1}+(n-1)d=21+(n-1)\cdot (-3)=-3n+24$.\\
		Khi  đó:
		\begin{itemize}
			\item  ${u_{11}}=-9$.
			\item $-6048=-3n+24\Rightarrow  n=2024$
		\end{itemize}
	} 
\end{ex}

% nguồn: Ok_dungsai_lop11_Lan2
\begin{ex}%[1D2H2-4]%[ma:fdcc52fc]
	
	Cho cấp số cộng $\left(u_n\right)$,  biết rằng $u_1=-3,u_6=27$, khi đó:
	\choiceTF 
		{Công sai của cấp số cộng bằng  $7$}
		 {\True Số hạng $u_{85}=501$}
		{Số hạng  $u_{10}=52$}
		{\True Tổng của 85 số hạng  đầu $S_{85}=21165$}
	\loigiai{ 
		Ta có $u_6=u_1+5d\Leftrightarrow 27=-3+5d\Leftrightarrow d=6$. Khi đó  $u_n=u_1+(n-1)d=-3+(n-1)\cdot 6=-9+6n$.
		\begin{itemchoice}
			\itemch  \textbf{Sai.}
			\itemch  \textbf{Đúng.}\\
			$u_{85}=-9+6\cdot85=501$. 
			\itemch \textbf{Sai.}\\ 
			$u_{10}=-9+6\cdot10=51$. 
			\itemch  \textbf{Đúng.}\\
			$S_{85}=\dfrac{85}{2}\left(2u_1+84d\right)=\dfrac{85}{2}[2\cdot (-3)+84\cdot 6]=21165$.. 
		\end{itemchoice}
	} 
\end{ex}

% nguồn: Ok_dungsai_lop11_Lan2
\begin{ex}%[1D2H2-4]%[ma:63f9cb81]
	
	Cho cấp số cộng $\left(u_n\right)$ có $u_5=18$ và $4S_n=S_{2n}$ (trong đó $S_n,S_{2n}$ theo thứ tự là tổng của $n$ và $2n$ số hạng đầu của  cấp số cộng).
	\choiceTF 
		{\True Số hạng đầu của cấp số cộng $\left(u_n\right)$ bằng $2$} 
		{Công sai của cấp số cộng $\left(u_n\right)$ bằng $3$} 
		{\True Số hạng  $u_{15}=58$}
		{Tổng $15$  số hạng đầu của cấp số cộng bằng $350$}
	 \loigiai{
		\begin{itemchoice}
			\itemch \textbf{Đúng.}\\ 
			Gọi $d$ là công sai của cấp số cộng, ta có $u_5=18\Leftrightarrow{u_1}+4d=18$.\\ 
			$4S_n=S_{2n}\Leftrightarrow\dfrac{4n}{2}\cdot \left[2u_1+(n-1)\cdot d\right]=\dfrac{2n}{2}\cdot \left[2u_1+(2n-1)\cdot  d\right]$\\
			$\Leftrightarrow 4u_1+(2n-2)\cdot d=2u_1+(2n-1) \cdot d \Leftrightarrow  2u_1-d=0$.\\
			Khi đó, ta  có $u_1=2,d=4$.
			\itemch \textbf{Sai.}\\ 
			Công sai của cấp  số cộng $\left(u_n\right)$ bằng $4$
			 \itemch \textbf{Đúng.}\\
			Số hạng tổng quát $u_n=2+(n-1)\cdot 4=4\cdot n-2$  suy ra $u_{15}=58$.
			\itemch  \textbf{Sai.}\\
			 Tổng $15$ số hạng đầu cấp số cộng là
			$S_{15}=\dfrac{15}{2}\left(2u_1+14d\right)=\dfrac{15}{2}(2\cdot 2+14\cdot  4)=450$.
		\end{itemchoice}
	} 
\end{ex}

% nguồn: Ok_dungsai_lop11_Lan2
\begin{ex}%[1D2H2-4]%[ma:0dffcec1]
	Cho cấp số cộng $(u_n)$ thoả mãn $\heva{& u_1-u_3+u_5=15\\& u_1+u_6=27}$. Xét tính  đúng sai của các mệnh đề sau
	 \choiceTF
		{\True  Số hạng $u_1=21$}
		{Công sai của cấp  số cộng bằng $-2$}
		{\True Số  hạng $u_{11}=-9$}
		{\True Số  $-6048$ là số hạng thứ $2024$}
	 \loigiai{
		Áp dụng công thức số hạng tổng quát của cấp  số cộng $u_n=u_1+(n-1)\cdot d$.\\
		Khi  đó
		\begin{align*}
			&\heva{&u_1-u_3+u_5=15\\& u_1+u_6=27}\\ 
			\Leftrightarrow &\heva{&  u_1-(u_1+2\cdot d)+(u_1+4\cdot d)=15\\&u_1+(u_1+5\cdot d)=27}\\
			\Leftrightarrow &\heva{& u_1+2\cdot d=15\\& 2u_1+5\cdot d=27}\\ 
			 \Leftrightarrow &\heva{& u_1=21\\& d=-3.}
		\end{align*}
		\begin{itemchoice}
			\itemch 
			Vì  $u_1=21$.
			\itemch 
			Vì  $d=-3$.
			 \itemch
			Vì  $u_n=u_1+(n-1)\cdot d=21+(n-1)\cdot (-3)=-3n+24$.\\
			Vậy $u_{11}=-3\cdot 11+24=-33+24=-9$. 
			\itemch 
			Vì ta  có $-6\,048=-3n+24 \Rightarrow -3n=-6\,072 \Rightarrow n=2\,024$.
		\end{itemchoice}
	 }
\end{ex}


%%%%%%%%%%%%%%%% TRẢ LỜI NGẮN %%%%%%%%%%%%%%%%

% nguồn: Ok_TraLoiNgan_Lop11_Lan2
\begin{ex}%[1D2H2-4]%[ma:a72d844f]
	Cho  cấp số cộng $(u_n)$ có $u_2=2\,001$ và $u_5=1\,995$. Khi đó $u_{1001}$ bằng
	\shortans{3} 
	 \loigiai{
		Ta có 
		$$\heva{&2\,001 = u_2 =  u_1+d\\&1\,995 = u_5 = u_1+4d} \Leftrightarrow \heva{&u_1 = 2\,003\\&d = -2} \Rightarrow u_{1001} = u_1+1\,000d = 3.$$
	 }
\end{ex}

% nguồn: Ok_TraLoiNgan_Lop11_Lan2
\begin{ex}%[1D2H2-4]%[ma:1b803aab]
	Câu 2.[Dự  án 11 HVA 2024-2025]
	Cho  cấp số cộng $(u_n)$ có ${u_2}=2023$ và $u_5=2032$. Khi đó $u_{15}$ bằng
	 \shortans{2062}
	\loigiai{ 
		Ta có $\heva{& {u_2}=2023 \\ & {u_5}=2032 }$ $\Leftrightarrow $ $\heva{& {u_1}+d=2023 \\ & {u_1}+4d=2032 }$ $\Leftrightarrow $ $\heva{& d=3 \\  & {u_1}=2020 }$ $\Rightarrow {u_{15}}=2020+14d=2062$.
	} 
\end{ex}

% nguồn: Ok_TraLoiNgan_Lop11_Lan2
\begin{ex}%[1D2H2-4]%[ma:d8c47828]
	 Câu 6[Dự án 11 HVA 2024-2025]
	Cho cấp số cộng $(u_n)$ có ${u_1}=-2;{u_3}=6$. Hỏi $2022$ là số hạng thứ bao nhiêu của cấp số  cộng đó?
	 \shortans{507}
	\loigiai{ 
		Ta có ${u_3}={u_1}+2d$ $\Rightarrow $ $6=-2+2d$  $\Rightarrow $ $ d=4$.\\
		Số hạng  tổng quát là ${u_n}={u_1}+(n-1)d=-2+(n-1)\cdot 4=4n-6$.\\
		Ta có $4n-6=2022$ $\Leftrightarrow $ $4n=2028$ $\Leftrightarrow $ $ n=507$.\\ 
	Vậy $2022$ là số hạng thứ $507$ của cấp số  cộng.}
\end{ex}

% nguồn: Ok_TraLoiNgan_Lop11_Lan2
\begin{ex}%[1D2H2-4]%[ma:296007d6]
	Cho cấp số cộng có tổng $n$ số hạng đầu là $S_n=3n^2+4n$, $n\in\mathbb{N}^*$. Giá trị của  số hạng thứ $10$ của cấp số cộng là
	\shortans{61} 
	\loigiai{ 
		Từ giả thiết ta có $S_1=u_1=3\cdot 1^2+4\cdot 1=7$.\\ 
		Ta có $S_n=3n^2+4n=\dfrac{n(8+6n)}{2}$ $=\dfrac{n(7+6n+1)}{2}$  $\Rightarrow u_n=6n+1$ $\Rightarrow {u_{10}}=61$.
	} 
\end{ex}

% nguồn: Ok_TraLoiNgan_Lop11_Lan2
\begin{ex}%[1D2H2-4]%[ma:c24f6464]
	Một cấp số cộng có số hạng đầu là $1$, công sai là $4$, tổng của $n$  số hạng đầu là $561$. Khi đó số hạng thứ $n$ của cấp số cộng đó là $u_n$ có giá trị là bao nhiêu?
	 \shortans{65}
	\loigiai{ 
		Ta có $\heva{&u_1=1,d=4\\&561=S_n=nu_1+\dfrac{n(n-1)}{2}d.}$\\ 
		Suy ra $561=n+\dfrac{n^2-n}{2}\cdot 4\Leftrightarrow 2n^2-n-561=0\Leftrightarrow n=17$.\\ 
		Vậy $u_n={u_{17}}=u_1+16d=1+16\cdot 4=65$. 
	 }
\end{ex}

% nguồn: Ok_TraLoiNgan_Lop11_Lan2
\begin{ex}%[1D2H2-4]%[ma:db8c9efc]
	Câu 3[Dự án  11 HVA 2024-2025]
	Một cấp số cộng có số hạng thứ  $5$ bằng $18$ và số hạng thứ $12$ bằng $32$. Tìm số hạng thứ $50$ của cấp số cộng này.
	 \shortans{108}
	 \loigiai{
		Ta biểu diễn số hạng thứ $5$ và số hạng thứ $12$ theo số hạng thứ nhất ${u_1}$ và công  sai $d$.\\
		Ta có $\heva{& {u_5}={u_1}+4d \\ & {u_{12}}={u_1}+11d }$ $\Leftrightarrow $ $\heva{& {u_1}+4d=18 \\ & {u_1}+11d=32 }$ $\Leftrightarrow  $ $\heva{& {u_1}=10 \\ & d=2 }$.\\
	Số hạng  thứ $50$ của cấp số cộng là ${u_{50}}={u_1}+49d=108 $.}
\end{ex}

% nguồn: Ok_TraLoiNgan_Lop11_Lan2
\begin{ex}%[1D2H2-4]%[ma:4ba6019b]
	 Câu 8[Dự án 11 HVA 2024-2025]
	Cho hai cấp số cộng $( {x_n})\colon 4;7;10\dots$ và $( {y_n})\colon 1;6;11\dots$. Hỏi trong $2018$  số hạng đầu tiên của mỗi cấp số có bao nhiêu số hạng chung?
	 \shortans{403}
	\loigiai{ 
		Số hạng tổng quát của cấp số cộng $( {x_n})$ là: ${x_n}=4+( n-1)\cdot 3$ $=3n+1$.\\ 
		Số hạng tổng  quát của cấp số cộng $( {y_n})$ là: ${y_m}=1+( m-1)\cdot 5$ $=5m-4$.\\
		Giả sử $ k$ là $1$ số hạng chung của hai cấp số cộng trong $2018$ số hạng đầu tiên của mỗi  cấp số.\\
		Vì $ k$ là $1$ số hạng của $( {x_n})$ nên $ k=3i+1$ với $1\le i\le 2018$ và  $ i\in {{\mathbb{N}}^*}$.\\
		Vì $ k$ là $1$ số hạng của $( {y_n})$ nên $ k=5j-4$ với $1\le j\le 2018$ và $ j\in {{\mathbb{N}}^*}$.\\ 
		Do đó $3i+1=5j-4$  $\Rightarrow 3i=5j-5 \Rightarrow i=\dfrac{5(j-1)}{3}$ $\Rightarrow i\vdots 5$ $\Rightarrow i\in \{ 5;10;15;\dots2015 \}$.\\
	Vậy  có $403$ số hạng chung.}
\end{ex}

% nguồn: Ok_TraLoiNgan_Lop11_Lan2
\begin{ex}%[1D2H2-4]%[ma:bc463264]
	Câu 4[Dự án 11 HVA 2024-2025] 
	Biết ${u_4}+{u_8}+{u_{12}}+{u_{16}}=224 $. Tính  ${u_{10}}$
	 \shortans{56}
	 \loigiai{
		Ta có ${u_4}+{u_8}+{u_{12}}+{u_{16}}=224$\\ 
		$\Leftrightarrow  $ $( {u_1}+3d)+( {u_1}+7d)+( {u_1}+11d)+( {u_1}+15d)=224$ $\Leftrightarrow $ $4{u_1}+36d=224$ $\Leftrightarrow $ ${u_1}+9d=56$.\\
	Khi đó ${u_{10}}={u_1}+9d=56$.} 
\end{ex}

% nguồn: Ok_TraLoiNgan_Lop11_Lan2
\begin{ex}%[1D2H2-4]%[ma:66ec5d27]
	 Câu 7[Dự án 11 HVA 2024-2025]
	Xác định công sai của cấp số cộng $( {u_n})$,  biết $\heva{& {u_9}-{u_4}=15 \\ & {u_3}\cdot {u_8}=184 }$.
	 \shortans{3}
	 \loigiai{Ta có
		$\heva{& {u_9}-{u_4}=15 \\ & {u_3}\cdot {u_8}=184 }$ $\Leftrightarrow $ $\heva{& {u_1}+8d-{u_1}-3d=15 \\ & ({u_1}+2d)\cdot({u_1}+7d)=184 }$ $\Leftrightarrow $ $\heva{& 5d=15 \\ & ({u_1}+2d)\cdot({u_1}+7d)=184 }$\\ 
	$\Leftrightarrow $ $\heva{& d=3 \\ & ({u_1}+6)\cdot({u_1}+21)=184 }$ $\Leftrightarrow $ $\heva{& d=3 \\ & u_1^2+27{u_1}-58=0 }$ $\Leftrightarrow  $ $\heva{& d=3 \\ & {u_1}=2;{u_1}=-29 }$.}
\end{ex}

% nguồn: Ok_TraLoiNgan_Lop11_Lan2
\begin{ex}%[1D2H2-4]%[ma:ae82486c]
	Cho cấp số cộng $(u_n)$ có ${d=-2~}$  và $S_8=72$ Tìm số hạng đầu tiên $u_1$.
	\shortans{16} 
	 \loigiai{
		 Ta có
		$\heva{&d=-2\\&72=8u_1+\dfrac{8\cdot 7}{2}d}\Rightarrow72=8u_1+28\cdot (-2)\Leftrightarrow u_1=16$. 
	} 
\end{ex}

% nguồn: Ok_TraLoiNgan_Lop11_Lan2
\begin{ex}%[1D2H2-4]%[ma:56d45d76]
	Một dãy số có các số hạng đầu tiên là $1$, $8$, $22$, $43$, ... Hiệu của hai số hạng liên tiếp của dãy số đó lập thành một cấp số cộng $7$, $14$, $21$, $...$, $7n$. Hỏi số $35\,351$ là  số hạng thứ mấy của dãy số đã cho?
	\shortans{101} 
	 \loigiai{
		Gọi $u_1,u_2,...,u_n$ lần lượt là số hạng thứ nhất, thứ hai, ...,  thứ $n$ của dãy số đã cho.\\
		 Ta có $\heva{&u_2-u_1=7\\&u_3-u_2=14\\&u_4-u_3=21\\& \cdots \\&u_n-{u_{n-1}}=7(n-1).}$\\
		 Cộng từng vế các phương trình của hệ ta được
		\[u_n-u_1=7+14+21+\cdots+7(n-1)=7\cdot  \dfrac{n(n-1)}{2}\quad(*).\]
		Thay $u_n=35\,351$ và $u_1=1$ vào ta  được
		\[35\,351-1=7\cdot \dfrac{n(n-1)}{2}\Leftrightarrow n^2-n-10\,100=0\Leftrightarrow \hoac{&n=101\\&n=-100.}\] 
		Với điều kiện $n\in \mathbb{N^*}$, ta  tìm được $n=101$.
	 }
\end{ex}

% nguồn: Ok_TraLoiNgan_Lop11_Lan2
\begin{ex}%[1D2H2-4]%[ma:07bdc208]
	Câu  5[Dự án 11 HVA 2024-2025]
	Cho cấp số cộng $( {u_n})$ xác định bởi  $\heva{& {u_1}=-2 \\ & {u_{n+1}}={u_n}+3 }$. Số $70$ là số hạng thứ bao nhiêu của cấp số cộng?
	 \shortans{25}
	 \loigiai{
		Ta có $\heva{& {u_1}=-2 \\ & {u_2}={u_1}+3=1  }$$\Rightarrow $ $ d={u_2}-{u_1}=3$\\
		Áp dụng công  thức số hạng tổng quát ${u_n}={u_1}+( n-1)d$.\\
		Ta có $70=-2+( n-1)\cdot 3$ $\Leftrightarrow  $ $70=-2+3n-3$ $\Leftrightarrow $ $n=25$.\\
	Vậy số $70$ là số hạng thứ $25$  của cấp số cộng.}
\end{ex}

% nguồn: Ok_TraLoiNgan_Lop11_Lan1
\begin{ex}%[1D2V2-4]%[ma:55991c65]
	Cho bốn số thực tạo thành một cấp số cộng có tổng bằng $28$ và tổng các bình phương của chúng bằng $276$. Tìm tích  của bốn số đó.
	\shortans{585} 
	\loigiai{ 
		 Gọi bốn số cần tìm theo thứ tự cấp số cộng là $a-3r$, $a-r$, $a+r$, $a+3r$.\\
		Ta  có $\heva{&a-3r+a-r+a+r+a+3r=28 \\
		 &(a-3r)^2+(a-r)^2+(a+r)^2+(a+3r)^2=276}$\\
		$\Leftrightarrow \heva{&4a=28 \\ 
		 &4a^2+20r^2=276}
		\Leftrightarrow \heva{&a=7 \\ 
		&r^2=4}\Leftrightarrow  \heva{&a=7\\
		&r=\pm 2.}$\\ 
		Vậy bốn số cần tìm  là $1$, $5$, $9$, $13$; tích của chúng bằng $585$.
	 }
\end{ex}

% nguồn: Ok_TraLoiNgan_Lop11_Lan1
\begin{ex}%[1D2V2-4]%[ma:8c26b0eb]
	Một cấp số cộng có $7$ số hạng. Biết rằng tổng của số hạng đầu và số hạng cuối bằng $30$, còn  tổng của số hạng thứ ba và số hạng thứ sáu bằng $35$. Tính số hạng thứ bảy của cấp số cộng đó.
	\shortans{30} 
	\loigiai{ 
		 Ta có: $\heva{&u_1+u_7=30\\&u_3+u_6=35}
		\Leftrightarrow  \heva{&2u_1+6d=30\\&2u_1+7d=35}
		\Leftrightarrow \heva{&u_1=0\\&d=5.}$\\ 
		Vậy  $u_7=u_1+6d=0+6\cdot5=30$.
	 }
\end{ex}

% nguồn: Ok_TraLoiNgan_Lop11_Lan1
\begin{ex}%[1D2V2-4]%[ma:0b0b495f]
	Cho phương trình $1+6+11+16+\cdots+x=970$. Biết rằng $1{,}6{,}11{,}\cdots x$ là  một cấp số cộng, Tìm $x$?
	\shortans{96} 
	 \loigiai{
		Vì  $1{,}6{,}11{,}\cdots x$ lập thành một cấp số cộng với $u_1 = 1$, công sai là $d = 6-1 =5$.\\
		Giả sử $x$ là số hạng thứ $n$ của cấp số  cộng nên ta có $S_n = 970$ hay là
		\[\dfrac{\left[2u_1 + \left(n-1\right)d \right]n  }{2} = 970 \Leftrightarrow \dfrac{\left[2 + \left(n-1\right)5 \right]n }{2} = 970 \Leftrightarrow 5n^2-3n-1940=0 \Leftrightarrow \hoac{&n=20\\&n=-\dfrac{97}{5}. }\]
		Suy ra $x$ là số hạng thứ $20$ nên $x = u_1 + 19d  = 1+19\cdot5 = 96$.
	} 
\end{ex}


%%%%%%%%%%%%%%%% TỰ LUẬN %%%%%%%%%%%%%%%%

% nguồn: Ok_TuanLuan_Toan11_Lan1
\begin{bt}%[1D2H2-4]%[1D2H2-6]%[ma:5c5616bd]
	\textbf{(1.0 điểm)} Cho cấp số  cộng $(u_n)$, biết
	$\heva{{u_2} - {u_3} + {u_5} = 12&\\{u_1} + {u_6} = 25}$. Hãy tìm số hạng đầu $u_1$, công sai $d$ và tổng $15$ số hạng  đầu tiên của cấp số cộng này.
	 \loigiai{
		\begin{itemize}
			\item $\heva{{u_2} - {u_3} + {u_5} = 12&\\{u_1} + {u_6} = 25}\Leftrightarrow $ $\heva{{{u_1} + d - \left(  {{u_1} + 2d} \right) + {u_1} + 4d = 12}&\\{{u_1} + {u_1} + 5d = 25}}\Leftrightarrow $
			 $\heva{{u_1} &= 15\\d &= - 1.}$
			\item 
			 ${S_{15}} = \dfrac{{15}}{2}\left[{2.15 + 14\left( { - 1} \right)} \right] = 120$.
		\end{itemize}
	 }
\end{bt}

% nguồn: Ok_TuanLuan_Toan11_Lan2
\begin{ex}%[1D2H2-4]%[1D2H2-6]%[ma:d6f870c7]
	 \textbf{(1.0 điểm)} Cho cấp số cộng $(u_n)$, biết
	$\heva{{u_2} - {u_3} + {u_5} = 12&\\{u_1} + {u_6} = 25}$. Hãy tìm  số hạng đầu $u_1$, công sai $d$ và tổng $15$ số hạng đầu tiên của cấp số cộng này.
	 \loigiai{
		\begin{itemize}
			\item $\heva{{u_2} - {u_3} + {u_5} = 12&\\{u_1} + {u_6} = 25}\Leftrightarrow $ $\heva{{{u_1} + d - \left( {{u_1} + 2d} \right) + {u_1} + 4d = 12}&\\{{u_1} + {u_1} +  5d = 25}}\Leftrightarrow $
			$\heva{{u_1}  &= 15\\d &= - 1.}$
			\item 
			${S_{15}}  = \dfrac{{15}}{2}\left[{2.15 + 14\left( { - 1} \right)} \right] = 120$.
		\end{itemize}
	} 
\end{ex}

% nguồn: Ok_TuanLuan_Toan11_Lan2
\begin{ex}%[1D2H2-4]%[ma:5c6010f6]
	Cho cấp số cộng $\left(u_{n}\right)$, biết $u_{1}=-5$, $d=2$. Số $81$ là số hạng thứ bao nhiêu của cấp số  cộng?
	\loigiai{Ta có $u_{n}=u_{1}+(n-1) d  \Leftrightarrow 81=-5+(n-1) 2 \Leftrightarrow n=44$.\\
	Vậy $81$ là số hạng  thứ $44$.}
\end{ex}

% nguồn: Ok_TuanLuan_Toan11_Lan2
\begin{ex}%[1D2H2-4]%[ma:86be0efd]
	Cho cấp số cộng $(u_n)$ có  $u_1+2\cdot u_5=0$ và $S_4=14$. \\
	Hãy  tính số hạng đầu $u_1$ và công sai $d$.
	\loigiai{ 
		 Ta có
		\begin{align*}
			u_5&=u_1+4d  \Rightarrow u_1+2 u_5=u_1+2(u_1+4d)=3u_1+8d=0. \tag{1}
		\end{align*}
		 Mặt khác:
		\begin{align*}
			S_4&=\dfrac{4}{2}\cdot(2u_1+3d)=2(2u_1+3d)=14\Rightarrow 2  u_1+3d=7. \tag{2}
		\end{align*}
		Từ (1) và (2), ta có hệ phương trình 
		 \[
		\heva{& 
		3u_1+8d=0\\& 
		 2u_1+3d=7}
		\Rightarrow 
		 \heva{&
		u_1=8\\& 
		d=-3.} 
		 \]
		Vậy  $u_1=8$ và $d=-3$.
	 }
\end{ex}

% nguồn: Ok_TuanLuan_Toan11_Lan2
\begin{ex}%[1D2H2-4]%[ma:47004982]
	Cho dãy số $\left(u_{n}\right)$ là một cấp số cộng có $u_{1}=3$ và công sai $d=4$. Biết tổng $n$ số hạng đầu  của dãy số $\left(u_{n}\right)$ là $S_{n}=253$. Tìm $n$.
	 \loigiai{Ta có
		\begin{eqnarray*}
			&  &S_{n}=\dfrac{n\left(2 u_{1}+(n-1) d\right)}{2}\\
			&\Leftrightarrow &  \dfrac{n(2\cdot 3+(n-1) \cdot 4)}{2}=253\\
			 &\Leftrightarrow &4 n^{2}+2 n-506=0\\
			 &\Leftrightarrow & \hoac{&n=11 \text{ (nhận)}\\ &n=-\dfrac{23}{2}\text{ (loại).}}
		\end{eqnarray*}
		Vậy  $n=11$.
	 }
\end{ex}

% nguồn: Ok_TuanLuan_Toan11_Lan2
\begin{ex}%[1D2V2-4]%[ma:ff19e013]
	Cho cấp số cộng có tổng $n$ số hạng đầu là $S_{n}=4 n^{2}+3 n$, $n \in \mathbb{N}^{*}$ thì số hạng thứ $10$ của cấp số cộng là  bao nhiêu?
	\loigiai{Theo công thức ta có $\dfrac{n\left(u_{1}+u_{n}\right)}{2}=4 n^{2}+3 n \Leftrightarrow u_{1}+u_{n}=8 n+6  \Rightarrow u_{n}=-u_{1}+8 n+6$.\\
	Mà $u_{1}=S_{1}=7$  do đó $u_{10}=-7+8\cdot 10+6=79$.}
\end{ex}

% nguồn: Ok_TuanLuan_Toan11_Lan2
\begin{ex}%[1D2V2-4]%[ma:ae26cfcf]
	Tìm $x$ sao cho $x+3; \,2x+1$ và ${5x+2}$ là ba số hạng liên tiếp  của một cấp số cộng.
	 \loigiai{
		Từ ${x+3;\, 2x+1}$ và ${5x+2}$ là ba số  hạng liên tiếp của một cấp số cộng ta suy ra
		 $${(x+3)+(5x+2)=2(2x+1) \Leftrightarrow x=-\dfrac{3}{2}}. $$
		Thử lại ta có ba số tìm được là ${\dfrac{3}{2},-2,-\dfrac{11}{2}}$ thoả  mãn bài toán. Vậy ${x=-\dfrac{3}{2}}$.
	} 
\end{ex}

% nguồn: Ok_TuanLuan_Toan11_Lan2
\begin{ex}%[1D2V2-4]%[ma:d2a5509e]
	Cho cấp số cộng  có tổng $n$ số hạng đầu là $S_{n}=3 n^{2}+4 n$, $n \in \mathbb{N}^{*}$. Giá trị của số hạng thứ $10$ của cấp số cộng là bao nhiêu?
	\loigiai{Từ giả thiết ta có $S_{1}=u_{1}=3\cdot 1^{2}+4\cdot  1=7$.\\
	Ta có $S_{n}=3 n^{2}+4 n=\dfrac{n(8+6 n)}{2}=\dfrac{n(7+6 n+1)}{2} \Rightarrow  u_{n}=6 n+1 \Rightarrow u_{10}=61$.}
\end{ex}
